\documentclass[11pt,a4paper]{article}
\usepackage{amsmath,amssymb,amsthm}
\usepackage[margin=2.5cm]{geometry}
\usepackage{hyperref}

\newtheorem{definition}{Definition}[section]
\newtheorem{theorem}[definition]{Theorem}
\newtheorem{proposition}[definition]{Proposition}
\newtheorem{remark}[definition]{Remark}
\newtheorem{conjecture}[definition]{Conjecture}

\title{Hyperseed Formalization:\\
  GPT o1 Does Not Know What It's Doing}
\author{ProtomegaTron (automated formalization)}
\date{2026-09-18}

\begin{document}
\maketitle

\begin{abstract}
We formalize Ben Goertzel's ``GPT o1 Does Not Know What It's Doing''
(2024-10-03) within the Hyperseed ontology. The article examines o1's
improvements and fundamental limitations through concrete programming
examples. Key mappings: pattern matching as base similarity without fiber
transport, lack of mental model as absent fiber, chain-of-thought as
shallow/performative fiber, brittle competence as thin fiber, and scaling
as base expansion without fiber improvement.
\end{abstract}

\tableofcontents

\section{Thin fiber: pattern matching without understanding}

\begin{theorem}[Pattern matching $\neq$ fiber transport --- claims 4--8, 18--19]
GPT o1 matches surface patterns (base similarity) without genuine fiber
transport (understanding why patterns hold and when they break):
\[
  \text{PatternMatch}(x) = \arg\max_y \text{sim}(x, y_{\text{train}})
  \neq \text{FiberTransport}(x, \text{understanding})
\]

Lacking a mental model of code execution means lacking fiber structure
over the code base. The system has base access (can read/write code tokens)
but no fiber (understanding of what the code means). Errors are diagnostic:
o1 fails in ways no genuine reasoner would (novel edge cases) while
succeeding in ways that impress pattern matchers (familiar configurations).
\end{theorem}

\section{Shallow fiber: chain-of-thought}

\begin{proposition}[CoT as performative fiber --- claims 9--12, 20]
Chain-of-thought adds a shallow layer of fiber (explicit reasoning steps)
but not deep fiber (genuine understanding). The fiber is performative:
\[
  F_{\text{CoT}} = \text{shallow}(\text{plausible traces})
  \neq F_{\text{deep}}(\text{backtracking, transfer, edge-case reasoning})
\]

CoT can simulate backtracking in text but can't genuinely reconsider
earlier steps --- it generates forward, always. This is performative fiber:
fiber that looks like reasoning in its surface form but lacks the structural
properties (reversibility, compositionality) of genuine reasoning fiber.
\end{proposition}

\section{Base expansion vs fiber improvement}

\begin{theorem}[Scaling limits --- claims 13--15, 21--22]
Scaling LLMs is base expansion, not fiber improvement:
\[
  \text{Scale} \uparrow \implies |B_{\text{patterns}}| \uparrow
  \quad \text{but} \quad \text{quality}(F) \not\uparrow
\]

Each generation shows diminishing returns on deep capabilities. The
architectural limit is fundamental: transformers don't support genuine
mental models, working memory, or systematic reasoning. These require
fiber-level changes (architectural innovation), not base-level changes
(more data/compute).
\end{theorem}

\section{Cross-references}

\begin{remark}[Connections]
\begin{itemize}
\item \textbf{OpenClaw: Amazing Hands for a Brain} (2026-02-03): same
  hands-vs-brain distinction applied to LLM scaling.
\item \textbf{Tensor Logic} (2025-12-16): tensor logic as genuine fiber
  improvement for reasoning, unlike CoT which is surface-level.
\item \textbf{Why Everyone Dies Gets AGI All Wrong} (2025-10-01): intelligence
  $\neq$ optimization parallels o1's pattern matching $\neq$ understanding.
\end{itemize}
\end{remark}

\end{document}
